\section*{Summary of findings}
\addcontentsline{toc}{section}{Summary of findings}

\paragraph{What was measured.} One linear program, three independent solvers.
The \FS{} of this project, a Dinkelbach iteration whose ratio oracle is a
minimum cut, and a divide-and-conquer parametric decomposition. All three return
the same object --- an optimal LP solution, its multipliers and the canonical
sequence of blocks --- and every result is re-checked by a separate process
before it is counted.

\paragraph{Correctness.} The three backends agree everywhere they were compared,
and agree with an exhaustive rational oracle on small instances and with an
independent LP solver on larger ones. Nothing in the campaigns contradicts the
correctness of any backend.

\generated{table_summary}

\paragraph{Speed.} The parametric backend is the only configuration that solves
the whole test set: $46$ of $46$. Dinkelbach solves $40$, the tuned \FS{} $36$
and the reference \FS{} $27$. Two families are solved by the parametric backend
alone --- disconnected graphs and the instances derived from MineLib --- where
the other three finish nothing within the limit.

The medians do not rescue the counts. Over the instances it solves, the tuned
\FS{} has a median of $0.142$\,s against the parametric backend's $0.057$\,s,
and that median is taken over the smaller subset it can finish; Dinkelbach sits
between them at $0.234$\,s but solves four more instances. On the historical
precedence-knapsack instances all four configurations except the reference
finish everything, and the ordering there is unambiguous: $0.027$\,s median for
the parametric backend, $0.087$\,s for Dinkelbach, $0.655$\,s for the tuned
\FS{}. The honest summary is that this implementation is neither the fastest
where it finishes nor able to finish everywhere.

\paragraph{What made the difference.} Not the pricing rule. The first
measurements had \FS{} one to three orders of magnitude behind, and parameter
search over pricing schemes moved the result by a few per cent. What moved it by
factors was the handling of degeneracy: the primal spends a median of $99.3\%$
of its pivots on zero steps, and three changes follow from that ---
releasing the anti-cycling fallback instead of letting it become the operating
rule, restricting the ratio-test scan to the arcs that can actually block, and
stopping that scan at the first zero-step blocker. The last two reproduce
techniques already present in the group's historical prototype.

The ablation confirms the ranking and corrects two of the search's choices.
Measured as paired speedups against the setting that forces Bland's rule,
the frozen configuration is $2.83$ times faster at the median and up to $18.8$
times on individual instances, so the anti-cycling release is worth a factor
rather than a percentage; moving the degeneracy trigger back to its old value
gives up most of that, $1.25$ instead of $2.83$. The ratio-test refinements are
worth about a factor of two each, $2.17$ and $2.20$.

Two settings the search did \emph{not} choose are faster than the one it did:
starting from the sink basis scores $3.33$ and enabling transitive reduction
$3.46$, against the frozen configuration's $2.83$. The frozen configuration is
reported as it was frozen, on training and validation data only, and not
silently improved after the fact --- but the ablation says plainly that the
search left something on the table, and \cref{sec:discussion} treats it as a
finding rather than a footnote.

\paragraph{How to read the tables.} Times are wall-clock medians over replicas of
the same instance and configuration. \emph{PAR-2} charges twice the time limit
for every instance a configuration did not solve, so it mixes speed and
robustness into one number; it is always reported next to the plain count of
solved instances so that a time-out is never hidden inside an average. Paired
speedups are computed only on the instances both settings solved and verified.
Performance profiles state the compared set and never count an unsolved instance
as covered.

\paragraph{What this report does not claim.} It does not claim novelty: the
audit of the literature was deliberately not carried out in this phase. It does
not claim anything about the state of the art in maximum flow: the flow backends
use a Dinic implementation written for this project, not a competitive flow code.
